\section{Conclusion}
\label{sec:conclusion}

We set out from the observation that Kolmogorov's axioms~\cite{kolmogorov1933} fix a logic
before they fix an arithmetic, and that nothing in the concept of
graded belief privileges the Boolean choice.  T
The entire classical surplus is one axiom, whose statical
face is the complement rule (Theorem~\ref{thm:hinge}) and whose
dynamical face is the agreement of Dempster and Bayesian
conditioning (Theorem~\ref{thm:cond-hinge}).  Everything else
survives the intuitionistic substrate intact: conditioning, Bayes
symmetry, disjoint additivity, total probability over genuine
partitions, and inclusion--exclusion itself
(Theorem~\ref{thm:incl-excl}).  The weaker logic adds the slack 
structure, along with its two-gap decomposition, and the
atomic/diffuse divide. For this new structure, 
we demonstrate a belief-function reading through
$\neg\neg$, and give a computability-theoretic reason the
non-degeneracy.

\paragraph{Open problems.}
Three are mathematical: (1) representing the diffuse part on
non-spatial frames (the localic Riesz direction),  (2) completing
the analytic residue of the Acz\'el development, namely continuity
of the generator and its extension past the cone to the rest of
the domain (the cone results and
the log-sum-exp example are done), and (3) a uniqueness
theorem: deriving modularity from a more primitive criterion,
knowing (Theorem~\ref{thm:no-functional}) that it cannot come from
disjunction data alone, with the mixture characterization
(Theorem~\ref{thm:mixture}) as the current best justification.  A
fourth is foundational hygiene at scale: a machine-checked
Dutch-book theorem in the style of~\cite{weatherson2003}, which to
our knowledge no proof assistant has for \emph{any} logic.

\paragraph{Functoriality and the monad question.}
Within a fixed doctrine, that is, a fixed choice of logic, the
connection to the categorical picture is concrete.  Pushforward along
a frame homomorphism $f : \Omega \to \Omega'$ sends a valuation $v$
on $\Omega'$ to $(f_* v)(a) = v(f(a))$ on $\Omega$, and this
preserves all four valuation axioms, making $\mathrm{Val}(-)$ a
functor on locales, the localic face of the functor part of the
Giry monad~\cite{giry1982}.  This is mechanized end to end
(\texttt{Valuation.comap} builds the pushforward, and
\texttt{comap\_id}/\texttt{comap\_comp} check the two functor laws
against identity and composition of frame homomorphisms).  The
point valuations $\delta_p$ of Section~\ref{sec:representation} are
the unit, already formalized, and Theorem~\ref{thm:finite-rep} shows
that on finite frames $\mathrm{Val}$ agrees, on objects, with the
classical finite-distribution monad on points.  The multiplication
is where the substance is.  Flattening a valuation of valuations is
integration, hence product valuations and Fubini, and by Kock's
identification of commutative monads with Fubini
theorems~\cite{kock2012}, whether statistics-as-kernel-composition
transports to a weaker logic is precisely the question whether that
logic's valuation functor is a commutative monad.  Two
machine-checked results delimit the multiplication from both sides.

On the chain product, the question of \emph{which} product valuation
to take does not arise, since the frame
$\mathrm{Low}(\mathbb{N} \times \mathbb{N})$ is \emph{rigid}: every
lower set is a \emph{finite} union of rectangles, because beyond a
stabilization height every occupied column is full, so the full
columns and finitely many rows account for everything
(\texttt{exists\_rect\_decomposition}).  The rigidity mechanism is
general.  By the inclusion--exclusion equality, two valuations
agreeing on any $\sqcap$-closed family agree on all finite joins of
its members (\texttt{eq\_on\_finsetSup\_of\_eq\_on}), and the chain
product is the instance where that family is the rectangles.  The
product valuation there, whenever it exists, is therefore unique,
with no continuity assumed (\texttt{Valuation.product\_unique}).

On the powerset frame $\mathrm{Set}(\mathbb{N} \times \mathbb{N})$,
whose elements are far from finite unions of rectangles,
non-uniqueness is a theorem.  The indicator of an ultrafilter is a
valuation (\texttt{ultrafilterValuation}), and any ultrafilter on a
product splits rectangles through its marginal pushforwards
(\texttt{rect\_mem\_iff}).  Taking a free ultrafilter $\mathcal{U}$
on $\mathbb{N}$, the two iteration orders of its tensor with itself
and its diagonal pushforward carry identical values on \emph{every}
rectangle, yet they differ on the upper triangle
$\{(m, m') \mid m < m'\}$
(\texttt{product\_valuation\_not\_\allowbreak unique}).  The two
orders of iterated integration of the same marginal data therefore
disagree: Fubini fails, concretely
(\texttt{fubini\_fails\_for\_valuations}).  Compare the Fubini
pathologies of non-additive
integration~\cite{denneberg1994,ghirardato1997}.  By Kock's
correspondence, no canonical commutative multiplication exists for
unrestricted valuations.  On the powerset frame these valuations are
finitely additive probability charges in the sense
of~\cite{bhaskara-rao1983}, so this half of the dichotomy is not
specific to intuitionism: the freedom is already present in the
classical finitely additive setting, and dropping continuity is what
admits it.  A monad therefore requires restricting either the
frames or the valuations.

\paragraph{The obstruction to a monad.}
The question upstream of the multiplication is worth stating
precisely, though what follows is an assessment rather than a
theorem.  A monad needs a category for $\mathrm{Val}$ to be an
endofunctor on.  The space of valuations must itself be an object, a
``space'' whose opens one can speak of.  For Scott-continuous
valuations that space is a locale and monad structure is available,
constructively, through the probabilistic powerdomain
line~\cite{jones-plotkin1989,vickers2008,coquand-spitters2009}.
Even there, exhibiting a cartesian closed category closed under the
construction is the Jung--Tix problem, open since 1998 for the
classical probabilistic powerdomain functor
itself~\cite{jung-tix1998}, though reframed versions of the category
have recently been settled.  Dropping continuity makes the
difficulty prior rather than merely worse, since the standard
presentations of the valuation space use continuity essentially.

One candidate route to a well-posed question does exist.  The
localic treatments present the space of continuous valuations as the
classifying locale of a propositional geometric theory~\cite{vickers2008,
coquand-spitters2009}, with generators asserting $q < v(a)$ for
rational $q$ and axioms for monotonicity, normalization, modularity
(expressible with rational witnesses and countable disjunctions),
and continuity.  Deleting the continuity axiom leaves a geometric
theory, so a classifying locale still exists on general grounds, and
classically its points are exactly the valuations of this paper.  We
have verified none of the structure this candidate would need.  The
unit looks unproblematic, but the multiplication requires
integrating the evaluation maps against a merely modular,
discontinuous valuation, where the sum of two lower-real-valued maps
is an infinite join, exactly where the diffuse part of
Theorem~\ref{thm:obstruction} leaks.  We offer the formulation as a
target, not a result.

A more modest strategy, with a better track record, is to carve out
a fragment rather than construct the whole monad.  On the
measure-theoretic side, the failure of cartesian closure for
measurable spaces was answered by quasi-Borel
spaces~\cite{heunen2017qbs}.  On the domain-theoretic side,
commutativity of the full valuations monad on \textsf{DCPO} is open,
and the response is a hierarchy of commutative \emph{submonads}
generated from simple valuations~\cite{jia2021commutative}.  Our
decomposition theorem names the corresponding fragment here: the
atomic part.  As a first machine-checked step in that direction, we
show that the countably-presented fragment composes correctly with
no space structure needed at all.  A countable mixture of valuations
is a valuation (\texttt{mixCountable}), a mixture of mixtures
flattens to a mixture with product weights, which is discrete
Chapman--Kolmogorov (\texttt{mixCountable\_mixCountable}), the
one-point mixture is the identity (\texttt{mixCountable\_unique}),
and iterated countable mixing is order-independent
(\texttt{mixCountable\_comm}), all by the unconditional Fubini of
$\ENNReal$ sums over sigma types.  This is associative, unital,
commutative kernel composition for countably-atomic data, in direct
contrast with the non-uniqueness above.  It is not yet a monad: the
presentations are not canonical, we do not quotient by ``presents the
same valuation,'' and the diffuse part is exactly what escapes every
countable presentation.

In summary, the unit, the finite level, and countable kernel
composition with its commutativity are proved.  Rigidity on chain
products and non-uniqueness on the powerset frame are proved, so the
general multiplication is settled: it is not canonically determined.
What remains is the ambient category, a formulation problem before
it is anything else.

\paragraph{Future work on the monad.}
Two continuations are concrete.  First, quotient the countable
presentations by ``presents the same valuation'' and package the
result as a Kleisli category of countably-atomic kernels, checking
that mixing descends to the quotient.  That would upgrade
composition-on-data to an honest categorical object.  Second,
mechanize the classifying-locale formulation sketched above, which
requires infrastructure \textsf{mathlib} currently lacks entirely:
classifying locales of propositional geometric theories.  Both
continuations inherit the boundary drawn here: whatever structure
exists must live on a fragment, because the unrestricted
multiplication is now provably not canonical.

Across doctrines, by contrast, we make no functoriality claim.
There is no ambient category in which ``logic $\mapsto$ its
probability theory'' is a functor, and each column requires
re-deriving what a state \emph{is} from the desiderata.  What is
uniform across columns is the method, and this paper's execution of
it is the Heyting column.

\paragraph{Future work beyond intuitionistic logic.}
Linear logic is a natural next column: its event algebras are
quantales, its conjunction is resource-sensitive, and a
``valuation'' there should grade the \emph{consumption} of evidence.
The analogue of Gleason's theorem or of our representation theorems,
that is, which states are compatible with the algebra, is open, as
is the right analogue of slack for the failure of contraction.
Existing probabilistic semantics of linear logic, probabilistic
coherence spaces~\cite{danos-ehrhard2011}, start from the additive
side rather than from the logic, so a Cox-style derivation over a
quantale remains to be done.  We expect the method used here, state
the classical desiderata, delete the axioms the logic no longer
licenses, and let a proof assistant referee what survives, what
splits, and what becomes unsatisfiable, to transfer directly.
